\chapter{Exploitation}
\label{chap:exploitation}

\section{Modèle d'exploitation}
L'infrastructure est exploitée par une seule personne au démarrage, avec une équipe qui se structure progressivement. L'exploitation doit donc être outillée pour minimiser la charge manuelle.

\section{Supervision}
\begin{tabularx}{\textwidth}{@{}lXl@{}}
\toprule
\textbf{Composant} & \textbf{Rôle} & \textbf{Statut} \\
\midrule
Prometheus & Métriques (CPU, RAM, disque, réseau) & À déployer \\
Grafana    & Tableaux de bord & À déployer \\
Loki       & Centralisation des logs & À déployer \\
Alertmanager & Alertes (email, Matrix, Teams) & À déployer \\
Uptime Kuma & Supervision de disponibilité des services & À déployer \\
\bottomrule
\end{tabularx}

\section{Sauvegardes}
\begin{tabularx}{\textwidth}{@{}lXX@{}}
\toprule
\textbf{Donnée} & \textbf{Fréquence} & \textbf{Rétention} \\
\midrule
Bases PostgreSQL (orchestration, OpenStack) & Quotidienne & 30 jours + mensuel 1 an \\
Volumes OpenStack & Quotidienne (incrémentale) & 30 jours \\
Configurations réseau (SRX, Cisco) & À chaque changement + hebdomadaire & 1 an \\
Images VM Proxmox & Hebdomadaire & 30 jours \\
Journal d'audit & Continu & 1 an (immuable) \\
\bottomrule
\end{tabularx}

\begin{okbox}[Règle 3-2-1]
Trois copies, deux supports, une copie off-site. La copie off-site peut être un bucket S3 compatible (Backblaze B2, Scaleway, OVH) ou un disque externe rotatif.
\end{okbox}

\section{Ticketing et gouvernance}
\begin{todobox}[À trancher]
Outils candidats : GLPI (mature, auto-hébergé), OTOBO, MantisBT, ou solution simple (Gitea + issues). Recommandation : GLPI pour la gestion des plaintes, demandes d'exclusion de groupe, escalades de quota, incidents.
\end{todobox}

\section{Runbooks prioritaires}
\begin{itemize}
  \item Création d'un étudiant (Keystone + projet + quotas).
  \item Création d'un groupe (projet + contribution initiale).
  \item Exécution d'une libération forcée (escalade ticket).
  \item Redémarrage d'un service OpenStack.
  \item Restauration d'un backup de base.
  \item Remplacement d'un disque défaillant.
  \item Réponse à une alerte SRX / Cisco.
\end{itemize}

\section{Astreinte et support}
\begin{itemize}
  \item Phase transitoire : best-effort, pas d'astreinte.
  \item Phase pilote : support aux heures ouvrées.
  \item Phase production : à définir avec la direction (support étendu ou non).
\end{itemize}

\section{Mises à jour}
\begin{itemize}
  \item Proxmox : mises à jour mensuelles, redémarrage planifié.
  \item OpenStack : une release majeure par an maximum, avec phase de test.
  \item Firmware : à chaque release constructeur critique (Dell, Juniper, Cisco).
  \item Politique de fenêtres de maintenance : annoncées 7 jours à l'avance.
\end{itemize}